\section{Marco teórico}
\label{sec:marco}
% Una subsección por concepto clave, solo los que se usan. Cada afirmación fuerte con referencia.

\subsection{Representación en punto flotante}
\label{subsec:punto-flotante}

Una computadora almacena cada número real en una cantidad fija de bits, así que solo puede
representar un subconjunto finito de $\mathbb{R}$. El estándar IEEE~754 fija cómo se elige ese
subconjunto y cómo se redondean los resultados de las operaciones (Institute of Electrical and
Electronics Engineers [IEEE], 2019). Python usa el
formato de doble precisión, \texttt{binary64}, que reparte sus 64 bits en un bit de signo, 11 bits
de exponente y 52 bits de mantisa almacenada. Todo número normal del formato se escribe como
\begin{equation}
  x = (-1)^{s}\,(1.m_1m_2\ldots m_{52})_2 \times 2^{e},
  \qquad -1022 \le e \le 1023 ,
  \label{eq:binary64}
\end{equation}
donde el 1 inicial no se guarda. La mantisa tiene entonces 53 bits de precisión efectiva. El mayor
valor representable es del orden de $1.80 \times 10^{308}$ y el menor normal positivo, del orden de
$2.23 \times 10^{-308}$.

Llamamos $\mathbb{F}$ al conjunto de números representables y $\mathrm{fl}(x)$ al elemento de
$\mathbb{F}$ más próximo a $x$, con desempate al de mantisa par, que es el modo de redondeo por
defecto del estándar. Dos constantes describen la resolución del formato. La unidad de redondeo
$u = 2^{-53} \approx 1.11 \times 10^{-16}$ acota el error relativo de representar un real,
\begin{equation}
  \frac{|\mathrm{fl}(x) - x|}{|x|} \le u ,
  \label{eq:redondeo}
\end{equation}
y el épsilon de máquina $\varepsilon_{\mathrm{M}} = 2^{-52} \approx 2.22 \times 10^{-16}$ mide la
distancia entre 1 y el siguiente elemento de $\mathbb{F}$. La distinción importa porque las dos
constantes difieren en un factor 2 y la bibliografía usa las dos (Higham, 2002).

De \eqref{eq:redondeo} sale el modelo estándar del análisis de error de redondeo: si
$\mathrm{op}$ es una de las cuatro operaciones aritméticas o la raíz cuadrada, el estándar exige que
el resultado en máquina sea el redondeo del resultado exacto, de modo que
\begin{equation}
  \mathrm{fl}(x \mathbin{\mathrm{op}} y) = (x \mathbin{\mathrm{op}} y)(1 + \delta),
  \qquad |\delta| \le u .
  \label{eq:modelo}
\end{equation}
Cada operación introduce un error relativo acotado por $u$, y el análisis de un algoritmo consiste
en seguir cómo se acumulan esos $\delta$ (Higham, 2002).

El espaciado de $\mathbb{F}$ no es uniforme: dentro de la binada $[2^{e}, 2^{e+1})$ los elementos
están separados por $2^{e-52}$. Los enteros son un caso particular útil. Todos los enteros con
$|n| \le 2^{53}$ pertenecen a $\mathbb{F}$; por encima de $2^{53}$ el espaciado pasa a valer 2 y
solo se representan los pares. De ahí sale el techo que aparece en el experimento de la
asociatividad. Si $x \ge 2^{54}$, el 1 queda por debajo de medio espaciado y el redondeo lo
descarta, de modo que $\mathrm{fl}(1 + x) = x$. En la binada $[2^{53}, 2^{54})$, donde el
espaciado vale 2, el 1 vale exactamente medio espaciado y el resultado lo decide el desempate a
mantisa par: $\mathrm{fl}(1 + x) = x$ si $x$ es múltiplo de 4 y $\mathrm{fl}(1 + x) = x + 2$ si no
lo es. El caso que interesa acá es el primero, porque $x = 2^{53}$ es múltiplo de 4.

El límite tampoco es simétrico, y esa asimetría explica una de las anomalías del experimento. En
$[2^{52}, 2^{53})$ el espaciado es 1 y en $[2^{53}, 2^{54})$ es 2, así que $2^{53} - 1$ es
representable exacto y $2^{53} + 1$ no lo es. Restar primero deja el resultado en la binada de
abajo, que tiene el doble de resolución, y ahí el 1 sobrevive: $\mathrm{fl}(1 - 2^{53})$ vale
$-(2^{53} - 1)$ y no $-2^{53}$.

Cuando una operación produce un valor de módulo mayor que el máximo representable, el estándar
devuelve $\pm\infty$ en lugar de detener el cálculo. Las operaciones indeterminadas sobre esos
valores, como $\infty - \infty$, devuelven \texttt{NaN}, un dato especial que se propaga por el
resto de la cuenta y que no es igual a sí mismo (IEEE, 2019).

\subsection{Enteros de precisión arbitraria}
\label{subsec:enteros}

El tipo \texttt{int} de Python no tiene ancho fijo: cada entero usa tantas palabras de memoria como
necesite su magnitud, y las operaciones aritméticas entre enteros son exactas mientras haya memoria
disponible (Python Software Foundation, 2024). No hay desborde ni redondeo, así que las identidades
algebraicas que valen en $\mathbb{Z}$ valen también en el programa.

Lo que sí cambia es el costo. Un entero de valor $|n|$ ocupa del orden de $\log_2 |n|$ bits, de modo
que $b^{k}$ necesita unos $k \log_2 b$ bits. Multiplicar ese número por $b$ cuesta un tiempo
proporcional a su longitud, y calcular las potencias sucesivas hasta $k = N$ acumula un costo del
orden de $N^{2}$. La aritmética exacta se paga en tiempo y en memoria.

C y Java no ofrecen este tipo por defecto: el tipo entero habitual tiene 32 o 64 bits fijos, y un
resultado fuera de rango desborda. Las dos familias resuelven el desborde de manera distinta: el
estándar de C declara que el desborde de un entero con signo es comportamiento indefinido, es decir,
que el compilador no está obligado a producir ningún resultado en particular (International
Organization for Standardization [ISO], 2018), mientras que Java define el resultado como la
reducción módulo $2^{n}$ en complemento a dos, con $n$ el ancho del tipo. Los enteros de ancho fijo
de NumPy, como \texttt{int64}, siguen esta segunda convención, que para ese tipo es la reducción
módulo $2^{64}$.

Esa convención importa para leer los datos. La
aritmética módulo $2^{64}$ es la de un anillo, con lo cual sigue siendo asociativa, conmutativa y
distributiva de forma exacta. Un cálculo cuyas operaciones intermedias desbordan puede terminar en
el resultado correcto si la expresión final vuelve a caer dentro del rango representable, porque las
cifras que se perdieron por arriba se cancelan entre sí.

\subsection{Asociatividad y conmutatividad en aritmética de punto flotante}
\label{subsec:asociatividad}

En $\mathbb{R}$ la suma es conmutativa y asociativa, y por eso una sumatoria se escribe sin
paréntesis y sin indicar en qué orden se recorren los sumandos. La suma de punto flotante conserva
solo la primera de las dos propiedades. Es conmutativa porque $\mathrm{fl}(x + y)$ depende del valor
exacto de $x + y$ y de la función de redondeo, y las dos son simétricas en sus argumentos. No es
asociativa: con $x = 1$ e $y = 2^{53}$,
\begin{equation}
  \mathrm{fl}\bigl(\mathrm{fl}(1 + 2^{53}) - 2^{53}\bigr) = 0
  \quad\text{y}\quad
  \mathrm{fl}\bigl(1 + \mathrm{fl}(2^{53} - 2^{53})\bigr) = 1 .
  \label{eq:contraejemplo}
\end{equation}
Los dos miembros son el mismo cálculo con distinta agrupación de paréntesis y difieren en todo el
valor del término (Goldberg, 1991).

La distinción tiene una consecuencia directa sobre el experimento del orden de suma. Una suma
acumulada de $N$ términos es una secuencia encadenada de sumas binarias, y permutar los sumandos
equivale a evaluar un agrupamiento distinto de la misma expresión. Lo que se pone a prueba al
reordenar los términos es la asociatividad y no la conmutatividad: el resultado cambia porque la
suma de punto flotante no es asociativa.

Para la suma acumulada de $N$ términos del mismo signo, el análisis de peor caso da una cota de
error absoluto proporcional a $N u \sum_{k} |x_k|$, que se obtiene suponiendo que los $N-1$
redondeos se acumulan todos en la misma dirección (Higham, 2002). Es una cota conservadora. Si en
cambio se modelan los $\delta$ de \eqref{eq:modelo} como variables independientes de media nula, la
suma de los errores se comporta como un paseo aleatorio y el crecimiento típico del error relativo
resulta
\begin{equation}
  E_{\mathrm{rel}} \sim \sqrt{N}\,u ,
  \label{eq:cota-raiz}
\end{equation}
que es la referencia estadística con la que se comparan los datos de este trabajo (Higham, 2002).
Del mismo análisis se desprende que, para términos positivos, acumular de menor a mayor módulo
reduce el error: cada término chico se suma cuando el acumulador todavía es chico y su información
no cae por debajo del último bit de la mantisa.

\subsection{Error relativo}
\label{subsec:error-relativo}

Dado un valor teórico $x \ne 0$ y su aproximación numérica $\hat{x}$, el error relativo es
\begin{equation}
  E_{\mathrm{rel}} = \frac{|\hat{x} - x|}{|x|} .
  \label{eq:error-relativo}
\end{equation}
Es la medida natural en punto flotante porque la precisión del formato es relativa: según
\eqref{eq:redondeo}, el error de representación crece con la magnitud del número, pero su cociente
con esa magnitud queda acotado por $u$ en todo el rango. Como referencia de lectura, un error
relativo del orden de $10^{-d}$ corresponde a unas $d$ cifras decimales significativas correctas.

La interpretación de esta medida tiene dos límites. El valor teórico también
se evalúa en punto flotante, así que el piso de lo que se puede medir es del orden de $u$: un error
relativo nulo significa que el resultado coincide bit a bit con el valor de referencia calculado, no
que sea exacto. El error relativo también necesita una forma cerrada contra la cual
comparar. Cuando la sucesión no tiene una, lo que se compara son dos representaciones entre sí, y la
diferencia absoluta $D_N$ entre ellas mide la discrepancia de los dos cálculos, no el error de
ninguno de los dos.

\subsection{Suma de Kahan}
\label{subsec:kahan}

La suma compensada de Kahan reduce el error de la suma acumulada arrastrando de forma explícita la
parte del término que el redondeo descarta (Kahan, 1965). El algoritmo mantiene, además del
acumulador $s$, una variable de compensación $c$ que guarda el error de la iteración anterior. En
cada paso, con el término $x_k$, calcula
\begin{equation}
  y = x_k - c, \qquad
  t = s + y, \qquad
  c = (t - s) - y, \qquad
  s = t ,
  \label{eq:kahan}
\end{equation}
donde cada línea se evalúa en punto flotante. La clave está en la tercera: cuando $|s| \ge |y|$, la
expresión $(t - s) - y$ recupera exactamente el error de redondeo cometido al formar $t = s + y$, y
la iteración siguiente lo devuelve al cálculo restándoselo al término entrante.

La cota de error del algoritmo para $N$ términos es de la forma $2u + O(Nu^{2})$ veces
$\sum_k |x_k|$, es decir, prácticamente independiente de $N$ para los tamaños de este trabajo
(Higham, 2002). Frente a la cota $Nu$ de la suma acumulada, la ganancia crece con la cantidad de
términos. El costo es de cuatro operaciones de punto flotante por término en lugar de una. Queda
además una advertencia de implementación: un compilador o una biblioteca que reasocie las
expresiones de \eqref{eq:kahan} suponiendo que la suma es asociativa anula la compensación, porque
en $\mathbb{R}$ la variable $c$ vale idénticamente cero.

\subsection{Sumas telescópicas}
\label{subsec:telescopicas}

Una suma es telescópica cuando su término general se escribe como diferencia de dos valores
consecutivos de una misma sucesión, $x_k = t_k - t_{k+1}$. En ese caso los términos intermedios se
cancelan de a pares y queda
\begin{equation}
  \sum_{k=1}^{N} (t_k - t_{k+1}) = t_1 - t_{N+1} .
  \label{eq:telescopica}
\end{equation}
La sucesión $b_N = \sum_{k=1}^{N} 1/(k(k+1))$ es de este tipo, con $t_k = 1/k$, porque la
descomposición en fracciones simples da $1/(k(k+1)) = 1/k - 1/(k+1)$. Aplicando
\eqref{eq:telescopica} se obtiene la forma cerrada
\begin{equation}
  b_N = 1 - \frac{1}{N+1} = \frac{N}{N+1} .
  \label{eq:forma-cerrada}
\end{equation}

Las cuatro expresiones anteriores son la misma en $\mathbb{R}$ y describen cuatro programas
distintos en $\mathbb{F}$. La suma del producto evalúa una división y una suma por término; la
telescópica evalúa dos divisiones y una resta antes de acumular; y las dos formas cerradas de
\eqref{eq:forma-cerrada} difieren en la cantidad de redondeos, una división y una resta contra una
sola división. Por \eqref{eq:modelo}, cada operación agrega su propio $\delta$, así que las cotas de
error de las cuatro no tienen por qué coincidir, y tampoco sus resultados.

\subsection{Cancelación catastrófica}
\label{subsec:cancelacion}

La cancelación catastrófica es la pérdida de cifras significativas que ocurre al restar dos números
de punto flotante próximos entre sí (Goldberg, 1991). Si los operandos llegan a la resta con errores
relativos $\delta_1$ y $\delta_2$ heredados de cálculos anteriores, el error relativo del resultado
queda acotado por
\begin{equation}
  \frac{|\widehat{x - y} - (x - y)|}{|x - y|}
  \le \frac{|x| + |y|}{|x - y|} \max(|\delta_1|, |\delta_2|) ,
  \label{eq:amplificacion}
\end{equation}
y el factor $(|x| + |y|)/|x - y|$ crece sin cota cuando $x$ e $y$ se acercan. Las cifras altas, que
eran correctas, se cancelan entre sí; las cifras bajas, que ya estaban contaminadas, quedan al
frente del resultado.

El nombre del fenómeno sugiere que el error lo introduce la resta, y no es así. El lema de Sterbenz
establece que si $x$ e $y$ son positivos y cumplen
$y/2 \le x \le 2y$, entonces $x - y$ pertenece a $\mathbb{F}$ y la resta se calcula sin error
(Sterbenz, 1974). Es decir que en el régimen donde la cancelación es más severa la resta suele ser
exacta: no crea error nuevo, expone el que traían los operandos.

La sucesión $c_N = \sum_{k=1}^{N} 1/(\sqrt{k^{2}+1} + k)$ ilustra el mecanismo. Su término general
admite la representación equivalente $\sqrt{k^{2}+1} - k$, que se obtiene multiplicando y dividiendo
por el conjugado. Para $k$ grande, $\sqrt{k^{2}+1} \approx k\,(1 + 1/(2k^{2}))$, de modo que los dos
operandos de la resta están dentro de un factor 2 y el lema de Sterbenz se aplica. El único error lo
aporta el redondeo previo de la raíz, de tamaño relativo $u$ y por lo tanto de tamaño absoluto del
orden de $k u$. Como el valor exacto de la diferencia es cercano a $1/(2k)$, el error relativo del
término calculado por la resta resulta del orden de
\begin{equation}
  \frac{k u}{1/(2k)} = 2k^{2} u ,
  \label{eq:cota-termino}
\end{equation}
una cota que crece cuadráticamente con $k$ y que alcanza el valor 1, es decir la pérdida completa de
las cifras significativas, en $k = 1/\sqrt{2u} = 2^{26}$. De ahí se desprende una cota para la suma
completa. Cada término aporta un error absoluto del orden de $k u$, de modo que si los errores de
los $N$ términos se alinearan en la misma dirección la discrepancia acumulada sería del orden de
$\sum_{k \le N} k u \approx N^{2} u / 2$, y crecería cuadráticamente con $N$. En lo que sigue se
usa la forma redondeada $N^{2}u$, que es el doble de esa estimación y por lo tanto la acota. Es el
análogo de la cota $Nu$ de la suma acumulada: describe el peor caso y no el típico. La representación como cociente suma en el
denominador en lugar de restar, así que no hay cancelación y su error relativo se mantiene del orden
de $u$ para todo $k$. La recomendación general es la de Goldberg (1991): frente a dos expresiones
algebraicamente equivalentes conviene elegir la que evita la resta de cantidades próximas, y la
racionalización por el conjugado es la transformación habitual para conseguirlo.
